% GENERATED FILE. Edit canonical lessons, then run npm run generate:books.
% canonical-source-hash: fnv1a64:d6c7a9b140ef01f7
% canonical-lessons: TA-C273-al, TA-C273-perror, TA-C273-utanpirantavar, TA-C273-irattaiyar, TA-C273-maittunan

\chapter{Person, Parents, Sibling, Twins, Brother-in-law}
\label{ch:ta-person-parents-sibling-twins-brother-in-law}
\hlchaptermodality{273}

\begin{chapteropening}
\textbf{By the end of this chapter:} \emph{I can say a person, parents, a sibling, twins and a brother-in-law.}

Complete the last lesson of chapter 273: I can say a person, parents, a sibling, twins and a brother-in-law.
\end{chapteropening}

\section[\texorpdfstring{\ta{ஆள்} (āḷ)}{ஆள் (āḷ)}]{\ta{ஆள்} (āḷ) — a person}

\label{lesson:TA-C273-al}



\begin{quote}
Before the new one: say the Tamil for a photograph, then the Tamil for a copy.
\end{quote}

\subsection*{You'll want to know: \ta{ஆள்}}
\textbf{\ta{ஆள்}} — \emph{āḷ} — \textquotedblleft{}a person\textquotedblright{}.

\begin{cousinweb}
\textbf{\ta{ஆள்}} (\emph{āḷ}) is native Dravidian. DEDR 399 lists Malayalam \emph{āḷ}, \textquotedblleft{}a person\textquotedblright{} and Kannada \emph{āḷ}.
\end{cousinweb}

\begin{grammarlens}[title={What you've built}]
One more word for this chapter.
\end{grammarlens}

\subsection*{Your turn}
\begin{itemize}
\raggedright
  \item \emph{Say it:} \emph{āḷ}
  \item \emph{Say it:} \emph{āḷ}, once more
  \item \emph{Say it:} say \emph{āḷ}
  \item \emph{Recall:} say the Tamil for a photograph, then the Tamil for a copy, then say \emph{āḷ} again
\end{itemize}

\begin{tcolorbox}[breakable,colback=teal!4,colframe=teal!35!black,title={Before you move on}]
What does \ta{ஆள்} mean? (\textquotedblleft{}A person\textquotedblright{}.) Say it once more.
\end{tcolorbox}
\section[\texorpdfstring{\ta{பெற்றோர்} (peṟṟōr)}{பெற்றோர் (peṟṟōr)}]{\ta{பெற்றோர்} (peṟṟōr) — parents}

\label{lesson:TA-C273-perror}



\begin{quote}
Before the new one: say the Tamil for a copy, then the Tamil for a person.
\end{quote}

\subsection*{You'll want to know: \ta{பெற்றோர்}}
\textbf{\ta{பெற்றோர்}} — \emph{peṟṟōr} — \textquotedblleft{}parents\textquotedblright{}.

\begin{cousinweb}
\textbf{\ta{பெற்றோர்}} (\emph{peṟṟōr}) is native: \textquotedblleft{}those who bore\textquotedblright{}, from \emph{peṟu}, \textquotedblleft{}to get, to bear (a child)\textquotedblright{} (DEDR 4422, which lists Malayalam peṟuka).
\end{cousinweb}

\begin{grammarlens}[title={What you've built}]
One more word for this chapter.
\end{grammarlens}

\subsection*{Your turn}
\begin{itemize}
\raggedright
  \item \emph{Say it:} \emph{peṟṟōr}
  \item \emph{Say it:} \emph{peṟṟōr}, once more
  \item \emph{Say it:} say \emph{peṟṟōr}
  \item \emph{Recall:} say the Tamil for a copy, then the Tamil for a person, then say \emph{peṟṟōr} again
\end{itemize}

\begin{tcolorbox}[breakable,colback=teal!4,colframe=teal!35!black,title={Before you move on}]
What does \ta{பெற்றோர்} mean? (\textquotedblleft{}Parents\textquotedblright{}.) Say it once more.
\end{tcolorbox}
\section[\texorpdfstring{\ta{உடன்பிறந்தவர்} (uṭaṉpiṟantavar)}{உடன்பிறந்தவர் (uṭaṉpiṟantavar)}]{\ta{உடன்பிறந்தவர்} (uṭaṉpiṟantavar) — a sibling}

\label{lesson:TA-C273-utanpirantavar}



\begin{quote}
Before the new one: say the Tamil for a person, then the Tamil for parents.
\end{quote}

\subsection*{You'll want to know: \ta{உடன்பிறந்தவர்}}
\textbf{\ta{உடன்பிறந்தவர்}} — \emph{uṭaṉpiṟantavar} — \textquotedblleft{}a sibling\textquotedblright{}.

\begin{cousinweb}
\textbf{\ta{உடன்பிறந்தவர்}} (\emph{uṭaṉpiṟantavar}) is native: \emph{uṭaṉ}, \textquotedblleft{}together with\textquotedblright{} (DEDR 945), and \emph{piṟantavar}, \textquotedblleft{}one who was born\textquotedblright{}, from \emph{piṟa}, \textquotedblleft{}to be born\textquotedblright{} (DEDR 4422): \textquotedblleft{}one born together with you\textquotedblright{}.
\end{cousinweb}

\begin{grammarlens}[title={What you've built}]
One more word for this chapter.
\end{grammarlens}

\subsection*{Your turn}
\begin{itemize}
\raggedright
  \item \emph{Say it:} \emph{uṭaṉpiṟantavar}
  \item \emph{Say it:} \emph{uṭaṉpiṟantavar}, once more
  \item \emph{Say it:} say \emph{uṭaṉpiṟantavar}
  \item \emph{Recall:} say the Tamil for a person, then the Tamil for parents, then say \emph{uṭaṉpiṟantavar} again
\end{itemize}

\begin{tcolorbox}[breakable,colback=teal!4,colframe=teal!35!black,title={Before you move on}]
What does \ta{உடன்பிறந்தவர்} mean? (\textquotedblleft{}A sibling\textquotedblright{}.) Say it once more.
\end{tcolorbox}
\section[\texorpdfstring{\ta{இரட்டையர்} (iraṭṭaiyar)}{இரட்டையர் (iraṭṭaiyar)}]{\ta{இரட்டையர்} (iraṭṭaiyar) — twins}

\label{lesson:TA-C273-irattaiyar}



\begin{quote}
Before the new one: say the Tamil for parents, then the Tamil for a sibling.
\end{quote}

\subsection*{You'll want to know: \ta{இரட்டையர்}}
\textbf{\ta{இரட்டையர்}} — \emph{iraṭṭaiyar} — \textquotedblleft{}twins\textquotedblright{}.

\begin{cousinweb}
\textbf{\ta{இரட்டையர்}} (\emph{iraṭṭaiyar}) is native Dravidian. DEDR 474 lists Malayalam \emph{raṇṭu}, \textquotedblleft{}two\textquotedblright{} and Telugu \emph{reṇḍu}. The entry lists \emph{iraṭṭai}, \textquotedblleft{}pair\textquotedblright{}, and \emph{iraṭṭaiyar}, \textquotedblleft{}twins\textquotedblright{}, under the word for \textquotedblleft{}two\textquotedblright{}.
\end{cousinweb}

\begin{grammarlens}[title={What you've built}]
One more word for this chapter.
\end{grammarlens}

\subsection*{Your turn}
\begin{itemize}
\raggedright
  \item \emph{Say it:} \emph{iraṭṭaiyar}
  \item \emph{Say it:} \emph{iraṭṭaiyar}, once more
  \item \emph{Say it:} say \emph{iraṭṭaiyar}
  \item \emph{Recall:} say the Tamil for parents, then the Tamil for a sibling, then say \emph{iraṭṭaiyar} again
\end{itemize}

\begin{tcolorbox}[breakable,colback=teal!4,colframe=teal!35!black,title={Before you move on}]
What does \ta{இரட்டையர்} mean? (\textquotedblleft{}Twins\textquotedblright{}.) Say it once more.
\end{tcolorbox}
\section[\texorpdfstring{\ta{மைத்துனன்} (maittuṉaṉ)}{மைத்துனன் (maittuṉaṉ)}]{\ta{மைத்துனன்} (maittuṉaṉ) — a brother-in-law}

\label{lesson:TA-C273-maittunan}



\begin{quote}
Before the new one: say the Tamil for a sibling, then the Tamil for twins.
\end{quote}

\subsection*{You'll want to know: \ta{மைத்துனன்}}
\textbf{\ta{மைத்துனன்}} — \emph{maittuṉaṉ} — \textquotedblleft{}a brother-in-law\textquotedblright{}.

\begin{cousinweb}
\textbf{\ta{மைத்துனன்}} (\emph{maittuṉaṉ}) is a loan from Sanskrit \emph{maithuna}, \textquotedblleft{}paired, coupled\textquotedblright{} (CDIAL 10341); DEDR App. 53 traces the kinship sense through Prakrit.
\end{cousinweb}

\begin{grammarlens}[title={What you've built}]
One more word for this chapter.
\end{grammarlens}

\subsection*{Your turn}
\begin{itemize}
\raggedright
  \item \emph{Say it:} \emph{maittuṉaṉ}
  \item \emph{Say it:} \emph{maittuṉaṉ}, once more
  \item \emph{Say it:} say \emph{maittuṉaṉ}
  \item \emph{Recall:} say the Tamil for a sibling, then the Tamil for twins, then say \emph{maittuṉaṉ} again
\end{itemize}

\begin{tcolorbox}[breakable,colback=teal!4,colframe=teal!35!black,title={Before you move on}]
What does \ta{மைத்துனன்} mean? (\textquotedblleft{}A brother-in-law\textquotedblright{}.) Say it once more.
\end{tcolorbox}
